\documentclass[11pt]{article}
\newcommand{\HomeworkHeader}{习题四}
% Shared LaTeX preamble for homework/*/main.tex

% --- Base layout ---
\usepackage[a4paper, top=2.5cm, bottom=2.5cm, left=2cm, right=2cm]{geometry}
\usepackage{fontspec}
\usepackage[UTF8]{ctex}%latex支持中文宏包
\usepackage{amsmath, amssymb, amsthm}
\usepackage{mathtools}
\usepackage[svgnames]{xcolor}
\usepackage{pgfplots}
\pgfplotsset{compat=1.18}
\usepackage[most]{tcolorbox}
\usepackage{enumitem}
% Modified: Change label from hyphen (-) to bullet point ($\bullet$) for better visibility
\setlist[itemize]{label=$\bullet$}
\setlist[enumerate]{label=(\arabic*)}


% --- Colors ---
\definecolor{primaryColor}{RGB}{46, 52, 64}
\definecolor{accentColor}{RGB}{94, 129, 172}
\definecolor{boxFill}{RGB}{236, 240, 241}
\definecolor{solLine}{RGB}{163, 190, 140}
\definecolor{expandColor}{RGB}{191, 97, 106}

% --- TikZ styles (DFA / NFA / PDA) ---
\usepackage{tikz}
\usetikzlibrary{automata, positioning, arrows.meta, bending, backgrounds}
\tikzset{
    dfa_state/.style={
        state,
        thick,
        draw=accentColor,
        fill=boxFill,
        text=primaryColor,
        minimum size=1.1cm
    },
    dfa_edge/.style={
        ->,
        >=stealth,
        thick,
        draw=primaryColor,
        auto
    },
    pda_node/.style={
        state,
        thick,
        draw=accentColor,
        fill=boxFill,
        text=primaryColor,
        minimum size=1.2cm
    },
    pda_edge/.style={
        ->,
        >=stealth,
        thick,
        draw=primaryColor,
        auto
    },
    rule_expand/.style={
        pda_edge,
        draw=expandColor,
        text=expandColor
    },
    rule_match/.style={
        pda_edge,
        draw=solLine,
        text=solLine!80!black
    }
}

% --- Header / footer ---
\usepackage{fancyhdr}
\pagestyle{fancy}
\setlength{\headheight}{13.6pt}%避免页眉高度警告
\fancyhf{}
\fancyhead[L]{\kaishu 中国科学院大学}
\fancyhead[C]{林俊杰2023K8009916012}
\fancyhead[R]{\kaishu 理论计算机}
\usepackage{lastpage}
\cfoot{\small 第 \thepage \ 页，共 \pageref{LastPage} 页}

% --- Problem box ---
\newtcolorbox[use counter=problemSeq]{problem}[1][]{
    enhanced,
    breakable,
    colback=boxFill,
    colframe=accentColor,
    coltitle=white,
    fonttitle=\bfseries\large,
    title={题目 ~\theproblemSeq \quad #1},
    boxrule=0.5mm,
    arc=3mm,
    drop shadow=black!15!white,
    attach boxed title to top left={xshift=0.5cm, yshift*=-3mm},
    boxed title style={colback=accentColor},
    % 关键修改：每当创建一个新问题盒子时，自动将步骤计数器重置为0
    code={\setcounter{stepcount}{0}}
}

% --- Solution 环境 (红色盒子样式，支持可选序号) ---
% 修改说明：
% 1. [1][] 表示接受一个可选参数，默认为空
% 2. title={...} 中使用 \ifstrempty 判断是否有参数，如果有则添加空格和参数
\newtcolorbox{solution}[1][]{
    enhanced, breakable,
    colback=red!3!white,        % 背景色：淡粉
    colframe=red!75!black,      % 边框色：深红
    coltitle=white,
    fonttitle=\bfseries\Large\itshape,
    title={Solution\ifstrempty{#1}{}{~#1}}, % 动态标题：Solution (1)
    boxrule=0.5mm,
    arc=3mm,
    drop shadow=black!15!white,
    attach boxed title to top left={xshift=0.5cm, yshift*=-3mm},
    boxed title style={colback=red!75!black},
    before skip=0.5em,          % 盒子前的间距
    after skip=2em,             % 盒子后的间距
    before upper={\setcounter{stepcount}{0}} % 每个解答开始时重置步骤计数
}

% --- Title ---
\newcommand{\makeCustomTitle}[1]{
    \begin{center}
        \vspace*{1cm}
        {\Huge \bfseries \color{primaryColor} 作业 #1}\\
        \vspace{0.5cm}
        {\Large \color{accentColor} \today}
        \vspace{1.2cm}
    \end{center}
}

% --- 自定义环境定义 ---

% --- 自定义计数器与环境 ---
\newcounter{stepcount}
\newcounter{casecount}[stepcount]
\newcounter{problemSeq} % 新增：专门用于 Problem 的独立计数器

% 【关键修改】挂载钩子：每当遇到 \section 命令（含 \section*），将 problemSeq 计数器重置为 0
\pretocmd{\section}{\setcounter{problemSeq}{0}}{}{}

% 2. 定义带自动编号的步骤环境
\newenvironment{proofstep}[1]
  {%
    \refstepcounter{stepcount}% 增加计数
    \par\vspace{1.5em}%
    \noindent\textbf{\large \thestepcount. #1}% 输出 "1. 标题"
    \par\vspace{0.5em}%
  }
  {\par}

% 3. 定义无编号的结论环境 (用于最后的结论，格式同步骤)
\newenvironment{proofresult}[1]
  {%
    \par\vspace{1.5em}%
    \noindent\textbf{\large #1}% 仅输出标题
    \par\vspace{0.5em}%
  }
  {\par}

% 4. 定义带自动编号的情形环境
\newenvironment{proofcase}[1]
  {%
    \refstepcounter{casecount}% 增加计数
    \par\vspace{1em}%
    \noindent\textbf{情形 \thecasecount：#1}% 输出 "情形 1：标题"
    \par\vspace{0.2em}%
  }
  {\par}
\begin{document}
\makeCustomTitle{习题四}
% 红框选题：6（全部两问）、9、10、13、17、18。
% 来源 assets/习题四.pdf：第1页 6；第2页 9、10、13；第3页 17、18。
% 对应教材第103--105页。

\begin{problem}[原题 6]
确定以下系统同时完全能控、完全能观时，参数的取值范围。
\begin{align*}
\text{（1）}\quad
&A=\begin{bmatrix}-1&1&a\\0&-2&1\\0&0&-3\end{bmatrix},\quad
B=\begin{bmatrix}0\\0\\1\end{bmatrix},\quad
C=\begin{bmatrix}0&0&1\end{bmatrix};\\[2mm]
\text{（2）}\quad
&A=\begin{bmatrix}0&0&1\\0&1&0\\-2&-3&-5\end{bmatrix},\quad
B=\begin{bmatrix}0\\1\\a\end{bmatrix},\quad
C=\begin{bmatrix}0&1&b\end{bmatrix}.
\end{align*}
\end{problem}
\begin{solution}
\begin{proofstep}{写出本题采用的两个秩判据}
对三阶单输入单输出系统，定义能控性矩阵与能观性矩阵
\[
\mathcal C=[B\ AB\ A^2B],\qquad
\mathcal O=\begin{bmatrix}C\\CA\\CA^2\end{bmatrix}.
\]
完全能控要求 $\operatorname{rank}\mathcal C=3$，完全能观要求
$\operatorname{rank}\mathcal O=3$。两矩阵都是三阶方阵，因此等价于
它们的行列式都非零。下面分别计算两问，最后取两个条件的交集。
\end{proofstep}

\begin{proofstep}{第（1）问：先检查能控性，再检查能观性}
因为 $B=e_3$，$AB$ 就是 $A$ 的第三列。再乘一次 $A$，得到
\[
AB=\begin{bmatrix}a\\1\\-3\end{bmatrix},\qquad
A^2B=\begin{bmatrix}-a+1-3a\\-2-3\\9\end{bmatrix}
=\begin{bmatrix}1-4a\\-5\\9\end{bmatrix}.
\]
故
\[
\mathcal C=\begin{bmatrix}0&a&1-4a\\0&1&-5\\1&-3&9\end{bmatrix},
\quad
\det\mathcal C=
\det\begin{bmatrix}a&1-4a\\1&-5\end{bmatrix}
=-5a-(1-4a)=-(a+1).
\]
因此完全能控当且仅当 $a\ne-1$。

再看能观性。$C=[0\ 0\ 1]$ 选取 $A$ 的第三行，故
$CA=[0\ 0\ {-3}]=-3C$，进一步 $CA^2=9C$。于是
\[
\mathcal O=\begin{bmatrix}0&0&1\\0&0&-3\\0&0&9\end{bmatrix},\qquad
\operatorname{rank}\mathcal O=1<3
\]
对任意 $a$ 都成立。其原因也可直接从方程看出：输出为 $y=x_3$，
而 $\dot x_3=-3x_3+u$ 与 $x_1,x_2$ 无关。在同一输入下改变 $x_1(0),x_2(0)$，
输出完全不变，因此无法恢复这两个初值。故两条件没有共同取值：
\[
\boxed{\text{不存在使第（1）个系统同时完全能控、完全能观的 }a.}
\]
\end{proofstep}

\begin{proofstep}{第（2）问：求完全能控的参数条件}
由矩阵乘法依次得到
\[
AB=\begin{bmatrix}a\\1\\-5a-3\end{bmatrix},\qquad
A^2B=\begin{bmatrix}-5a-3\\1\\-2a-3+25a+15\end{bmatrix}
=\begin{bmatrix}-5a-3\\1\\23a+12\end{bmatrix}.
\]
于是
\[
\mathcal C=\begin{bmatrix}
0&a&-5a-3\\
1&1&1\\
a&-5a-3&23a+12
\end{bmatrix}.
\]
按第一行展开，其行列式为
\begin{align*}
\det\mathcal C
&=-a\bigl((23a+12)-a\bigr)
  +(-5a-3)\bigl((-5a-3)-a\bigr)\\
&=-a(22a+12)+(5a+3)(6a+3)\\
&=8a^2+21a+9.
\end{align*}
方程 $8a^2+21a+9=0$ 的判别式为
$21^2-4\cdot8\cdot9=153=9\cdot17$，故两根为
\[
a=\frac{-21\pm3\sqrt{17}}{16}.
\]
除去这两个值，系统完全能控。
\end{proofstep}

\begin{proofstep}{第（2）问：求完全能观的参数条件并取交集}
分别计算输出行与 $A$ 的乘积：
\[
CA=[0\ 1\ b]A=[-2b\ \ 1-3b\ \ -5b],
\]
\[
CA^2=[-2b\ \ 1-3b\ \ -5b]A
=[10b\ \ 1+12b\ \ 23b].
\]
因此
\[
\mathcal O=\begin{bmatrix}
0&1&b\\
-2b&1-3b&-5b\\
10b&1+12b&23b
\end{bmatrix}.
\]
按第一行展开并合并同类项：
\begin{align*}
\det\mathcal O
&=-\bigl((-2b)(23b)-(-5b)(10b)\bigr)\\
&\quad+b\bigl((-2b)(1+12b)-(1-3b)(10b)\bigr)\\
&=-4b^2+b(6b^2-12b)=2b^2(3b-8).
\end{align*}
所以完全能观当且仅当 $b\ne0$ 且 $b\ne\tfrac83$。
能控条件只含 $a$，能观条件只含 $b$，二者同时成立的完整范围为
\[
\boxed{a\notin\left\{\frac{-21+3\sqrt{17}}{16},\frac{-21-3\sqrt{17}}{16}\right\},
\qquad b\notin\left\{0,\frac83\right\}.}
\]
\end{proofstep}
\end{solution}

\begin{problem}[原题 9]
设 $n$ 阶单输入单输出系统 $(A,b,c)$ 既能控又能观。证明：当 $n\ge2$ 时，$A$ 与 $bc$ 不可交换。
\end{problem}
\begin{solution}
\begin{proofstep}{把可交换假设写成可以利用的等式}
本题 $b$ 是 $n\times1$ 列向量、$c$ 是 $1\times n$ 行向量，因此 $bc$ 是 $n\times n$
矩阵，可以讨论它与 $A$ 是否交换。
若 $b=0$，系统必不能控；若 $c=0$，系统必不能观。故题设保证 $b,c$ 均非零。

反设 $A$ 与 $bc$ 可交换，则由矩阵乘法的结合律有
\[
(Ab)c=b(cA).
\]
这个等式把 $Ab$ 与 $b$ 联系起来。为了去掉左边的行向量 $c$，取
\[
v=\frac{c^{\mathsf T}}{cc^{\mathsf T}}.
\]
由于 $c\ne0$，分母 $cc^{\mathsf T}>0$，且 $cv=1$。
\end{proofstep}

\begin{proofstep}{证明输入方向必须成为一个特征方向}
将交换等式右乘 $v$，得到
\[
(Ab)(cv)=b(cAv),\qquad
Ab=\lambda b,\quad \lambda=cAv.
\]
$cAv$ 为标量，所以这表明 $b$ 是 $A$ 的右特征向量。
继续乘 $A$，有 $A^2b=A(\lambda b)=\lambda Ab=\lambda^2b$；归纳得
$A^kb=\lambda^kb$。也就是说，输入方向及其经系统矩阵反复传播所得的方向都平行于 $b$。
\end{proofstep}

\begin{proofstep}{用能控矩阵的秩得到矛盾}
于是能控性矩阵成为
\[
\mathcal C=[b\ Ab\ \cdots\ A^{n-1}b]
=b[1\ \lambda\ \cdots\ \lambda^{n-1}].
\]
因为 $b\ne0$，所有列均在一维空间 $\operatorname{span}\{b\}$ 中，且第一列非零，所以
\[
\operatorname{rank}\mathcal C=1<n\qquad(n\ge2).
\]
这与完全能控所要求的秩 $n$ 矛盾。因此原假设不成立，即
\[
\boxed{A(bc)\ne(bc)A.}
\]
$n\ge2$ 的条件不可省略：当 $n=1$ 时二者均为标量，自然可交换。
\end{proofstep}
\end{solution}

\begin{problem}[原题 10]
$n$ 阶单输入单输出系统 $(A,b,c)$ 满足
\[
cA^{n-1}b=a\ne0,\qquad cA^kb=0\quad(k=0,1,\ldots,n-2).
\]
证明该系统既能控又能观。
\end{problem}
\begin{solution}
\begin{proofstep}{将题设的标量乘积组织成一个方阵}
要同时证明能控和能观，需要证明下列两个 $n\times n$ 矩阵可逆：
\[
\mathcal C=[b\ Ab\ \cdots\ A^{n-1}b],\qquad
\mathcal O=\begin{bmatrix}c\\cA\\\vdots\\cA^{n-1}\end{bmatrix}.
\]
题设给出的量恰好是二者相乘后的元素，因此考虑
\[
H=\mathcal O\mathcal C
=\begin{bmatrix}
cb&cAb&\cdots&cA^{n-1}b\\
cAb&cA^2b&\cdots&cA^nb\\
\vdots&\vdots&&\vdots\\
cA^{n-1}b&cA^nb&\cdots&cA^{2n-2}b
\end{bmatrix}.
\]
用从0开始的行列编号，$H_{ij}=(cA^i)(A^jb)=cA^{i+j}b$。
当 $i+j<n-1$ 时，题设保证该元素为0；当 $i+j=n-1$ 时，该元素为 $a$。
所以反对角线全为非零数 $a$，它左上侧全为0，右下侧元素无需知道。
\end{proofstep}

\begin{proofstep}{逆置列顺序，将已知结构变为三角形}
令 $J$ 为逆序置换矩阵，即 $J_{i,n-1-i}=1$、其余元素为0。
右乘 $J$ 相当于把第1列与第 $n$ 列的位置互换，并依此逆置全部列。
新矩阵的元素为
\[
(HJ)_{ij}=cA^{i+n-1-j}b.
\]
当 $i<j$ 时指数小于 $n-1$，元素为0；当 $i=j$ 时指数为 $n-1$，元素为 $a$。
因此 $HJ$ 是下三角矩阵，主对角线全为 $a$，其行列式为
\[
\det(HJ)=a^n\ne0.
\]
列序完全逆置共有 $n(n-1)/2$ 个逆序对，故
$\det J=(-1)^{n(n-1)/2}$。由 $\det(HJ)=\det H\det J$，得到
\[
\det H=(-1)^{n(n-1)/2}a^n\ne0.
\]
\end{proofstep}

\begin{proofstep}{从乘积可逆推出两个秩判据同时成立}
因为 $H=\mathcal O\mathcal C$，有
\[
\det\mathcal O\,\det\mathcal C=\det H\ne0.
\]
两个数的乘积非零，说明它们分别非零。因此 $\mathcal C$ 可逆、秩为 $n$，系统完全能控；
$\mathcal O$ 也可逆、秩为 $n$，系统完全能观。
本证明不需要 $cA^kb$ 在 $k\ge n$ 时的具体取值，且对 $n=1$ 也成立。
\end{proofstep}
\end{solution}

\begin{problem}[原题 13]
定出下列单输入单输出系统的能观规范型和变换矩阵：
\[
\dot x=\begin{bmatrix}-1&-2&-2\\0&-1&1\\1&0&1\end{bmatrix}x
+\begin{bmatrix}2\\0\\1\end{bmatrix}u,
\qquad y=\begin{bmatrix}1&1&0\end{bmatrix}x.
\]
\end{problem}
\begin{solution}
\begin{proofstep}{确认系统完全能观}
将题面矩阵记为 $A,b,c$。左乘输出行 $c=[1\ 1\ 0]$，相当于取 $A$ 前两行之和，故
\[
cA=\begin{bmatrix}-1&-3&-1\end{bmatrix},\qquad
cA^2=\begin{bmatrix}0&5&-2\end{bmatrix},
\]
\[
\mathcal O=\begin{bmatrix}c\\cA\\cA^2\end{bmatrix}
=\begin{bmatrix}1&1&0\\-1&-3&-1\\0&5&-2\end{bmatrix}.
\]
按第一行展开，得
\[
\det\mathcal O=1\bigl(6+5\bigr)-1\bigl(2-0\bigr)=9\ne0.
\]
因此系统完全能观，可以构造三阶能观规范型。
\end{proofstep}

\begin{proofstep}{用特征多项式确定目标规范型}
计算特征多项式：
\begin{align*}
\det(sI-A)
&=\det\begin{bmatrix}s+1&2&2\\0&s+1&-1\\-1&0&s-1\end{bmatrix}\\
&=(s+1)^2(s-1)+2+2(s+1)\\
&=s^3+s^2+s+3.
\end{align*}
采用教材第4.4节（第95--97页）的末列伴随型，目标矩阵与输出行为
\[
A_o=\begin{bmatrix}0&0&-3\\1&0&-1\\0&1&-1\end{bmatrix},\qquad
c_o=[0\ 0\ 1].
\]
最后一列依次是特征多项式常数项、一次项、二次项系数的相反数。
选择 $c_o$ 后，第三个新状态直接等于输出，即 $z_3=y$。
\end{proofstep}

\begin{proofstep}{从输出开始构造新坐标的三行}
先写 $z=Sx$，将 $S$ 的三行记为 $\ell_1,\ell_2,\ell_3$。
要使 $y=z_3$，必须取 $\ell_3=c$。
要使新状态矩阵为 $A_o$，必须满足 $SA=A_oS$，也就是
\[
\ell_1A=-3\ell_3,\qquad
\ell_2A=\ell_1-\ell_3,\qquad
\ell_3A=\ell_2-\ell_3.
\]
从最后一个等式往前计算，得到
\[
\ell_2=cA+c,\qquad
\ell_1=\ell_2A+c=cA^2+cA+c.
\]
第一个等式由 Cayley--Hamilton 定理保证：
\[
\ell_1A=c(A^3+A^2+A)=-3c=-3\ell_3.
\]
所以从原坐标到新坐标的变换为
\[
S=\begin{bmatrix}cA^2+cA+c\\cA+c\\c\end{bmatrix}
=\begin{bmatrix}0&3&-3\\0&-2&-1\\1&1&0\end{bmatrix},\qquad \det S=-9.
\]
这里 $\det S\ne0$，所以新坐标没有丢失任何状态信息。
若把坐标变换写成 $x=Tz$，则 $T=S^{-1}$，由消元得到
\[
\boxed{T=S^{-1}=\begin{bmatrix}
-\tfrac19&\tfrac13&1\\
\tfrac19&-\tfrac13&0\\
-\tfrac29&-\tfrac13&0
\end{bmatrix},\qquad x=Tz.}
\]
这里 $S$ 将旧坐标映为新坐标，$T$ 将新坐标映回旧坐标，二者方向相反。
\end{proofstep}

\begin{proofstep}{计算输入列和输出行，写出完整规范型}
由于 $z=Sx$，有
\[
\dot z=SAS^{-1}z+Sbu,\qquad y=cS^{-1}z.
\]
第三步已保证 $SAS^{-1}=A_o$。另外，直接计算
\[
Sb=\begin{bmatrix}0&3&-3\\0&-2&-1\\1&1&0\end{bmatrix}
\begin{bmatrix}2\\0\\1\end{bmatrix}
=\begin{bmatrix}-3\\-1\\2\end{bmatrix},
\]
而 $S$ 的第三行是 $c$，所以 $cS^{-1}=e_3^{\mathsf T}SS^{-1}=[0\ 0\ 1]$。
最终得到
\[
\boxed{
\dot z=\underbrace{\begin{bmatrix}0&0&-3\\1&0&-1\\0&1&-1\end{bmatrix}}_{SAS^{-1}}z
+\underbrace{\begin{bmatrix}-3\\-1\\2\end{bmatrix}}_{Sb}u,
\qquad y=\underbrace{\begin{bmatrix}0&0&1\end{bmatrix}}_{cS^{-1}}z.}
\]
\end{proofstep}

\begin{proofstep}{核对规范型的输入输出关系}
从各行方程出发，在零初始条件下有
\[
sZ_1=-3Z_3-3U,\qquad
sZ_2=Z_1-Z_3-U,\qquad
sZ_3=Z_2-Z_3+2U.
\]
由最后一式 $Z_2=(s+1)Z_3-2U$；代入第二式，得到
$Z_1=(s^2+s+1)Z_3-(2s-1)U$。
再代入第一式，得
\[
(s^3+s^2+s+3)Z_3=(2s^2-s-3)U.
\]
由于 $Y=Z_3$，传递函数为
\[
G(s)=\frac{2s^2-s-3}{s^3+s^2+s+3},
\]
与直接计算原系统的 $c(sI-A)^{-1}b$ 相同。
输入列各分量也可以通过教材公式核对：
\[
\beta_2=cb=2,\quad
\beta_1=cAb+cb=-1,\quad
\beta_0=cA^2b+cAb+cb=-3
\]
恰好按 $[\beta_0\ \beta_1\ \beta_2]^{\mathsf T}$ 的顺序组成 $Sb$。
\end{proofstep}
\end{solution}

\begin{problem}[原题 17]
线性定常系统为 $\dot x=Ax+Bu$、$y=Cx$。采用输出反馈 $u=Fy+v$ 后，闭环系统为
\[
\dot x=(A+BFC)x+Bv,\qquad y=Cx.
\]
证明 $(A+BFC,B,C)$ 能控能观的充要条件为 $(A,B,C)$ 能控能观。
\end{problem}
\begin{solution}
\begin{proofstep}{明确反馈符号与需要比较的判据}
设 $B\in\mathbb R^{n\times m}$、$C\in\mathbb R^{p\times n}$、$F\in\mathbb R^{m\times p}$。
按题面的正号反馈 $u=Fy+v=FCx+v$，代入原方程得到
$\dot x=(A+BFC)x+Bv$。闭环的外部输入变为 $v$，输入矩阵仍为 $B$，输出矩阵仍为 $C$。

记 $A_f=A+BFC$。PBH 判据要求：对任意 $\lambda\in\mathbb C$，
\[
(A_f,B)\text{ 能控}\iff
\operatorname{rank}[\lambda I-A_f\ B]=n\quad\text{对所有 }\lambda;
\]
\[
(A_f,C)\text{ 能观}\iff
\operatorname{rank}\begin{bmatrix}\lambda I-A_f\\C\end{bmatrix}=n
\quad\text{对所有 }\lambda.
\]
因此只需证明反馈前后这两个矩阵的秩分别相同。
\end{proofstep}

\begin{proofstep}{比较能控性：构造可逆的列变换}
在能控 PBH 矩阵中，反馈多出的 $-BFC$ 是输入列 $B$ 的线性组合。
因此可以用列变换消去它。具体地，令
\[
R_F=\begin{bmatrix}I_n&0\\-FC&I_m\end{bmatrix},\qquad
R_F^{-1}=\begin{bmatrix}I_n&0\\FC&I_m\end{bmatrix}.
\]
逐块相乘得到
\[
[\lambda I-A-BFC\ \ B]
=[\lambda I-A\ \ B]
R_F.
\]
第一组列为 $(\lambda I-A)I_n+B(-FC)$，第二组列为
$(\lambda I-A)0+BI_m$，正是左边的两组列。
由于右乘可逆矩阵不改变秩，故
\[
\operatorname{rank}[\lambda I-A-BFC\ \ B]
=\operatorname{rank}[\lambda I-A\ \ B].
\]
对所有 $\lambda$ 同时使用此式，得到
$(A+BFC,B)$ 能控当且仅当 $(A,B)$ 能控。
\end{proofstep}

\begin{proofstep}{比较能观性：构造可逆的行变换}
能观 PBH 矩阵中，多出的 $-BFC$ 是输出行 $C$ 的线性组合。
相应地取行变换矩阵
\[
L_F=\begin{bmatrix}I_n&-BF\\0&I_p\end{bmatrix},\qquad
L_F^{-1}=\begin{bmatrix}I_n&BF\\0&I_p\end{bmatrix}.
\]
直接相乘，上方行块为 $I_n(\lambda I-A)-BFC$，下方行块为 $C$，所以
\[
\begin{bmatrix}\lambda I-A-BFC\\C\end{bmatrix}
=L_F\begin{bmatrix}\lambda I-A\\C\end{bmatrix}.
\]
左乘可逆矩阵不改变秩，故对于每个 $\lambda$，反馈前后的能观 PBH 秩均相同。
于是 $(A+BFC,C)$ 能观当且仅当 $(A,C)$ 能观。
\end{proofstep}

\begin{proofstep}{合并能控、能观两项结论}
若原系统能控且能观，第二、三步分别保证闭环也能控且能观，这是充分性。
若闭环能控且能观，利用同样的秩等式反向推回，原系统也能控且能观，这是必要性。
综上，
\[
\boxed{(A+BFC,B,C)\text{ 能控能观}\iff(A,B,C)\text{ 能控能观}.}
\]
整个论证只使用可逆行、列变换，没有对 $F$ 的取值作限制，也没有假定闭环稳定。
\end{proofstep}
\end{solution}

\begin{problem}[原题 18]
开环系统为
\[
A=\begin{bmatrix}0&1\\-1&0\end{bmatrix},\qquad
B=\begin{bmatrix}0\\1\end{bmatrix},\qquad
C=\begin{bmatrix}1&0\end{bmatrix}.
\]
证明该系统能控能观，但不存在常数输出反馈增益 $K$ 使 $A+BKC$ 的特征根全在左半平面。
\end{problem}
\begin{solution}
\begin{proofstep}{检查开环系统的能控性与能观性}
先计算
\[
AB=\begin{bmatrix}0&1\\-1&0\end{bmatrix}\begin{bmatrix}0\\1\end{bmatrix}
=\begin{bmatrix}1\\0\end{bmatrix},\qquad
CA=[1\ 0]\begin{bmatrix}0&1\\-1&0\end{bmatrix}=[0\ 1].
\]
因此
\[
\mathcal C=[B\ AB]=\begin{bmatrix}0&1\\1&0\end{bmatrix},\qquad
\mathcal O=\begin{bmatrix}C\\CA\end{bmatrix}
=\begin{bmatrix}1&0\\0&1\end{bmatrix}.
\]
其行列式分别为 $-1$ 和 $1$，都不为零，所以系统完全能控且完全能观。
\end{proofstep}

\begin{proofstep}{写出所有允许的静态输出反馈闭环}
本题输入、输出都是标量，故静态输出反馈矩阵 $K$ 实际上只能是一个实数 $k$。
由 $y=x_1$，正号反馈 $u=ky+v$ 等于 $u=kx_1+v$。
因此
\[
BkC=\begin{bmatrix}0\\1\end{bmatrix}k[1\ 0]
=\begin{bmatrix}0&0\\k&0\end{bmatrix},
\]
\[
A_k=A+BkC=\begin{bmatrix}0&1\\k-1&0\end{bmatrix},\qquad
\det(sI-A_k)=\det\begin{bmatrix}s&-1\\1-k&s\end{bmatrix}=s^2+1-k.
\]
任意 $k$ 都只能改变特征多项式的常数项，一次项始终为零。
\end{proofstep}

\begin{proofstep}{证明不存在两个实部均为负的闭环特征值}
因为 $\operatorname{tr}A_k=0$，两特征值满足 $\lambda_1+\lambda_2=0$。
若它们实部都严格小于0，那么
$\operatorname{Re}(\lambda_1+\lambda_2)<0$，与和为0矛盾。
故仅由迹就已证明：不存在使 $A_k$ 为 Hurwitz 矩阵的 $k$。

为了看清不同参数下具体发生什么，再解 $s^2+1-k=0$：
\begin{itemize}
\item $k<1$ 时，特征值为 $\pm\mathrm i\sqrt{1-k}$，均在虚轴上且互异。
自由运动是无阻尼振荡，虽有界却不趋于零，因而不是渐近稳定。
\item $k=1$ 时，$A_k=\begin{bmatrix}0&1\\0&0\end{bmatrix}$，有 $A_k^2=0$。
此时 $e^{A_kt}=I+tA_k$，故 $x_1(t)=x_1(0)+tx_2(0)$、$x_2(t)=x_2(0)$。
只要 $x_2(0)\ne0$ 就会线性增长，系统不稳定。
\item $k>1$ 时，特征值为 $\pm\sqrt{k-1}$，其中一个严格为正，沿相应特征方向指数增长，系统不稳定。
\end{itemize}
因此
\[
\boxed{\nexists\,k\in\mathbb R\quad\text{使 }A+BkC\text{ 为 Hurwitz 矩阵}.}
\]
\end{proofstep}

\begin{proofstep}{说明结论为何不与能控、能观性矛盾}
能控性保证可以通过完整状态反馈配置极点，但这里使用的静态输出反馈只允许访问
$y=x_1$，无法直接加入 $x_2$ 的反馈项。
例如，若允许完整状态反馈 $u=-x_1-2x_2+v$，则闭环矩阵变为
\[
\begin{bmatrix}0&1\\-2&-2\end{bmatrix},\qquad
\det(sI-A_{\rm cl})=s^2+2s+2,
\]
其特征值为 $-1\pm\mathrm i$，确实渐近稳定。
但 $u=-x_1-2x_2+v$ 不能写成本题允许的 $u=ky+v$。
同样，能观性意味着可以利用一段时间内的输入输出信息重构状态，并不意味着
单个时刻的标量 $y$ 已等于完整状态。因此这些性质与本题结论相容。
\end{proofstep}
\end{solution}
\end{document}
